\documentclass[12pt,a4paper]{article}
\usepackage[utf8]{inputenc}
\usepackage[T5]{fontenc}
\usepackage{lmodern} % Fix lỗi hiển thị font tiếng Việt (chống răng cưa/mất nét)
\usepackage[vietnamese]{babel}
\usepackage{geometry}
\geometry{a4paper, top=2cm, bottom=2cm, left=2.5cm, right=2cm}
\usepackage{hyperref}
\usepackage{xcolor}
\usepackage{tcolorbox}
\tcbuselibrary{skins,breakable}
\usepackage{enumitem}

% Thiết lập màu sắc
\definecolor{primary}{RGB}{33, 85, 163}
\definecolor{darkgray}{RGB}{60, 60, 60}

% Thiết lập layout box cho mỗi bài báo
\newtcolorbox{paperbox}[1]{
    colback=white,
    colframe=primary,
    colbacktitle=primary,
    title={\textbf{#1}},
    fonttitle=\bfseries\large,
    arc=2mm,
    boxrule=1pt,
    breakable,
    left=4mm, right=4mm, top=3mm, bottom=3mm,
    title filled=true,
    toptitle=2mm, bottomtitle=2mm
}

% Định dạng cho các tiêu đề mục con
\newcommand{\subheading}[1]{\vspace{0.3cm}\noindent\textbf{\textcolor{primary}{\MakeUppercase{#1}:}}\par\vspace{0.1cm}}

% Cài đặt danh sách
\setlist[itemize]{leftmargin=*, label=\textcolor{primary}{$\bullet$}, itemsep=2pt, parsep=2pt}

\begin{document}

\begin{center}
    {\huge\bfseries\color{primary} TỔNG HỢP NGHIÊN CỨU KHOA HỌC}\\[0.5cm]
    {\Large\bfseries CHỦ ĐỀ: NHẬN DIỆN BIỂN BÁO GIAO THÔNG (TRAFFIC SIGN RECOGNITION)}\\[1cm]
\end{center}

% --- BÀI SỐ 1 ---
\begin{paperbox}{1. Toward a new dataset: Mexican Traffic Signs ReWaIn-MTS for detection using deep learning}
    \textbf{Nguồn:} \url{https://ieeexplore.ieee.org/document/11045643}\\

    \subheading{Vấn đề giải quyết}
    \begin{itemize}
        \item Xây dựng và huấn luyện mô hình deep learning để nhận diện biển báo giao thông Mexico với hiệu suất cao hơn nhờ tăng cường dữ liệu.
    \end{itemize}

    \subheading{Phương pháp}
    \begin{itemize}
        \item Sử dụng các mô hình object detection hiện đại như YOLOv5, YOLOv8 và DETR.
        \item Đề xuất một chiến lược data augmentation nhằm cải thiện hiệu suất huấn luyện trên một dataset nhỏ dành riêng cho Mexico.
        \item Các mô hình được huấn luyện lại với dữ liệu tăng cường và đánh giá bằng mAP.
    \end{itemize}

    \subheading{Ưu điểm}
    \begin{itemize}
        \item Sử dụng các mô hình state-of-the-art (YOLOv8, DETR).
        \item Có dataset riêng cho Mexico mang tính thực tế.
        \item Áp dụng data augmentation giúp cải thiện mAP.
        \item So sánh nhiều mô hình khác quan.
    \end{itemize}

    \subheading{Nhược điểm}
    \begin{itemize}
        \item Dataset khá nhỏ (2,283 ảnh) và chỉ thu thập ở 8 thành phố, chưa đủ đa dạng.
        \item Không đề cập đến tốc độ xử lý (FPS) hoặc khả năng chạy real-time trên thiết bị nhúng.
        \item Không rõ performance trong điều kiện ánh sáng phức tạp hoặc bị che khuất.
    \end{itemize}

    \subheading{Kết quả \& Kết luận}
    \begin{itemize}
        \item Các mô hình được huấn luyện với tăng cường dữ liệu đạt hiệu suất (mAP) cao hơn đáng kể các phương pháp tiêu chuẩn.
        \item Việc áp dụng tăng cường dữ liệu đóng vai trò quan trọng và giúp cải thiện đáng kể hiệu suất phát hiện và phân loại với dataset nhỏ.
    \end{itemize}
\end{paperbox}

\vspace{0.5cm}

% --- BÀI SỐ 2 ---
\begin{paperbox}{2. A Robust Intelligent System for Text-Based Traffic Signs Detection and Recognition in Challenging Weather Conditions}
    \textbf{Nguồn:} \url{https://ieeexplore.ieee.org/document/10530242}\\
    
    \subheading{Vấn đề giải quyết}
    \begin{itemize}
        \item Phát hiện và nhận dạng chữ trên biển báo giao thông trong nhiều điều kiện thời tiết khác nhau.
    \end{itemize}

    \subheading{Phương pháp}
    \begin{itemize}
        \item Đề xuất semi-pipeline system gồm 3 giai đoạn:
        \item Dùng YOLOv5s tùy chỉnh để phát hiện bảng biển.
        \item Áp dụng thuật toán MSER để định vị vùng văn bản.
        \item Sử dụng OCR kết hợp NLP để nhận dạng và xử lý văn bản.
    \end{itemize}

    \subheading{Ưu điểm}
    \begin{itemize}
        \item Giải quyết vấn đề mới mẻ: text-based traffic sign.
        \item Hệ thống nhiều tầng chuyên biệt cho từng nhiệm vụ.
        \item Dùng YOLOv5s nhẹ có khả năng chạy nhanh và hệ thống hoạt động được trong nhiều điều kiện thời tiết với bộ ASAYAR dataset.
    \end{itemize}

    \subheading{Nhược điểm}
    \begin{itemize}
        \item Hệ thống nhiều bước dễ gặp lỗi dây chuyền (YOLO sai $\rightarrow$ MSER sai $\rightarrow$ OCR sai).
        \item MSER là thuật toán truyền thống nên có thể kém robust hơn Deep Learning và OCR phụ thuộc vào chất lượng ảnh.
        \item Không phải là pipeline end-to-end Deep Learning hoàn toàn.
    \end{itemize}

    \subheading{Kết quả}
    \begin{itemize}
        \item Hệ thống đạt hiệu suất khả quan trên MSTAR và ASAYAR dataset, có khả năng cạnh tranh với các phương pháp hiện đại (SOTA). Tuy nhiên abstract không cung cấp số liệu cụ thể.
    \end{itemize}
\end{paperbox}

\vspace{0.5cm}

% --- BÀI SỐ 3 ---
\begin{paperbox}{3. Deep Learning for Large-Scale Traffic-Sign Detection and Recognition}
    \textbf{Nguồn:} \url{https://ieeexplore.ieee.org/abstract/document/8709983}\\

    \subheading{Vấn đề giải quyết}
    \begin{itemize}
        \item Xây dựng hệ thống có thể xử lý và nhận dạng một số lượng rât lớn các loại biển báo giao thông (lên đến 200 loại) để phục vụ việc tự động hóa quản lý kho biển báo (traffic-sign inventory management).
    \end{itemize}

    \subheading{Phương pháp}
    \begin{itemize}
        \item Sử dụng mô hình chính Mask R-CNN theo cách end-to-end learning (toàn bộ pipeline được huấn luyện cùng lúc).
        \item Đề xuất một số cải tiến và áp dụng trên dataset mới gồm 200 loại biển báo.
    \end{itemize}

    \subheading{Ưu điểm}
    \begin{itemize}
        \item Hệ thống hoạt động theo hướng end-to-end, xử lý được category lớn.
        \item Có phân tích toàn diện về độ biến thiên (variation) trong cùng một loại biển.
        \item Phù hợp triển khai hệ thống quản lý vào thực tế.
    \end{itemize}

    \subheading{Nhược điểm}
    \begin{itemize}
        \item Chưa tối ưu về tốc độ xử lý nhanh và đòi hỏi lượng tài nguyên tính toán lớn.
        \item Còn thiếu so sánh chi tiết với các phương pháp phát hiện đối tượng SOTA khác.
    \end{itemize}

    \subheading{Kết quả \& Kết luận}
    \begin{itemize}
        \item Tỷ lệ lỗi (Error rate) dưới 3\% $\rightarrow$ Độ chính xác trên 97\% và áp dụng được cho 200 loại ngay cả khi độ biến thiên hình ảnh cao.
        \item Hiệu suất đạt mức đủ để triển khai trong ứng dụng thực tế.
    \end{itemize}
\end{paperbox}

\vspace{0.5cm}

% --- BÀI SỐ 4 ---
\begin{paperbox}{4. Robust Stacking Ensemble Model for Traffic Sign Detection and Recognition}
    \textbf{Nguồn:} \url{https://ieeexplore.ieee.org/document/10763484}\\

    \subheading{Vấn đề giải quyết}
    \begin{itemize}
        \item Xây dựng hệ thống phát hiện và nhận dạng biển báo giao thông dạng chữ trong môi trường thực tế khó khăn (ánh sáng yếu, sương mù, nhiễu, che khuất, chữ nhỏ/nghiêng). Mục tiêu hệ thống đủ nhanh để đáp ứng thời gian thực cho ITS.
    \end{itemize}

    \subheading{Phương pháp}
    \begin{itemize}
        \item Chuyển đổi không gian màu LUV (dùng kênh L tăng tương phản).
        \item Dùng MSER lọc bằng đặc trưng hình học để phát hiện chữ.
        \item Dùng YOLOv5s (với Mosaic, GIoU Loss) để phát hiện biển báo.
        \item OCR + NLP để nhận diện và hiệu chỉnh ký tự.
    \end{itemize}

    \subheading{Ưu điểm}
    \begin{itemize}
        \item Kết hợp Deep Learning và truyền thống, có tiền xử lý giúp giải quyết cực tốt ánh sáng kém.
        \item YOLOv5 cho tốc độ cao, giảm false positive.
        \item Định hướng tốt hệ thống real-time cho ITS.
    \end{itemize}

    \subheading{Nhược điểm}
    \begin{itemize}
        \item Phụ thuộc chất lượng ảnh đầu vào, độ tương phản thấp có thể làm MSER thất bại.
        \item Chưa phân tích chi tiết FPS và test khả năng chạy thiết bị nhúng.
    \end{itemize}

    \subheading{Kết quả \& Kết luận}
    \begin{itemize}
        \item Độ chính xác cao, giảm false positive, hoạt động hiệu quả thời gian thực. Cải thiện đáng kể độ tin cậy trong môi trường Giao thông Thông minh (ITS).
    \end{itemize}
\end{paperbox}

\vspace{0.5cm}

% --- BÀI SỐ 5 ---
\begin{paperbox}{5. Traffic Sign Classification for Autonomous Vehicles Using Split and Federated Learning Underlying 5G}
    \textbf{Nguồn:} \url{https://ieeexplore.ieee.org/document/10288388}\\
    \textit{(Lưu ý: Nội dung tóm tắt tương tự nội dung bài 1. Nhưng mã số DOI IEEE khác)}

    \subheading{Vấn đề \& Đóng góp chính}
    \begin{itemize}
        \item Giải quyết bài toán thiếu hụt mẫu dữ liệu bằng Data Augmentation chuyên sâu cho biển báo.
        \item Khắc phục hạn chế của việt huấn luyện dataset nhỏ (2,283 ảnh tại 8 thành phố Mexico).
    \end{itemize}

    \subheading{Phương pháp}
    \begin{itemize}
        \item Triển khai YOLOv5, YOLOv8 và DETR với chiến lược làm giàu dữ liệu đặc biệt.
        \item Đánh giá hiệu suất mAP qua từng điều kiện augmentation.
    \end{itemize}

    \subheading{Ưu điểm \& Nhược điểm}
    \begin{itemize}
        \item \textbf{Ưu:} State-of-the-art architectures (YOLOv8, DETR), dataset thực tế Mexico, so sánh đối chiếu khách quan.
        \item \textbf{Nhược:} Kích thước dữ liệu quá nhỏ bé để tổng quát hoá. Thiếu phân tích FPS và tính chạy môi trường thực tiễn (edge devices).
    \end{itemize}

    \subheading{Kết quả}
    \begin{itemize}
        \item Data Augmentation đóng vai trò quan trọng không thể thiếu, tăng kết quả mAP lên đáng kể so với việc train tiêu chuẩn.
    \end{itemize}
\end{paperbox}

\vspace{0.5cm}

% --- BÀI SỐ 6 ---
\begin{paperbox}{6. Real-time traffic sign recognition based on a general purpose GPU and deep-learning}
    \textbf{Nguồn:} \url{https://journals.plos.org/plosone/article?id=10.1371/journal.pone.0173317}\\

    \subheading{Vấn đề giải quyết}
    \begin{itemize}
        \item Hệ thống nhận dạng bị giảm độ chính xác và tốc độ khi di chuyển trong ánh sáng yếu hoặc thay đổi ánh sáng phức tạp.
    \end{itemize}

    \subheading{Phương pháp}
    \begin{itemize}
        \item Mô hình phân cấp (hierarchical model): sử dụng thuật toán Byte-MCT để chống biến đổi sáng, AdaBoost phát hiện vùng biển báo, SVM xác minh vùng, CNN phân loại biển báo.
        \item Tăng tốc quy trình bằng việc tính toán xử lý song song trên GPGPU.
    \end{itemize}

    \subheading{Ưu điểm}
    \begin{itemize}
        \item Vô cùng mạnh mẽ trước ánh sáng bất lợi. Nhờ GPGPU hệ thống đạt tốc độ real-time tuyệt vời.
    \end{itemize}

    \subheading{Nhược điểm}
    \begin{itemize}
        \item Lệ thuộc nặng vào thiết bị phần cứng cao cấp (phải có GPU chuyên).
        \item Cấu trúc tuần tự nhiều bước (MCT $\rightarrow$ AdaBoost $\rightarrow$ SVM $\rightarrow$ CNN) phức tạp và tốn công finetune từng đoạn.
    \end{itemize}

    \subheading{Kết quả}
    \begin{itemize}
        \item F1-score đạt $0.97$ trên tập dữ liệu tổng hợp (Đức \& Hàn Quốc). Ổn định tốt cả vào ban đêm.
    \end{itemize}
\end{paperbox}

\vspace{0.5cm}

% --- BÀI SỐ 7 ---
\begin{paperbox}{7. Evaluation Method of Deep Learning-Based Embedded Systems for Traffic Sign Detection}
    \textbf{Nguồn:} \url{https://ieeexplore.ieee.org/abstract/document/9490209}\\

    \subheading{Vấn đề giải quyết}
    \begin{itemize}
        \item Phát hiện biển báo giao thông (Traffic Sign Detection - TSD) là một tác vụ cực kỳ quan trọng đối với xe tự hành, nhưng thường xuyên bị ảnh hưởng bởi các điều kiện bên ngoài (thời tiết, ánh sáng, phông nền phức tạp) và giới hạn xử lý bên trong.
        \item Dù đã có nhiều thuật toán Deep Learning, nhưng hiện tại đang thiếu một phương pháp luận và khung thực nghiệm chuẩn (dựa trên phân tích thống kê) để giúp các nhà nghiên cứu/kỹ sư đưa ra quyết định chính xác khi lựa chọn kết hợp giữa thuật toán và phần cứng (như CPU, GPU, TPU hay hệ thống nhúng). Bài báo ra đời nhằm lấp đầy khoảng trống này.
    \end{itemize}

    \subheading{Phương pháp}
    \begin{itemize}
        \item \textbf{Mô hình triển khai:} Tác giả kết hợp các mô hình phân loại (MobileNet v1 và ResNet50 v1) với thuật toán phát hiện đối tượng SSD (Single Shot Multibox Detector) và mạng đặc trưng FPN (Feature Pyramid Network).
        \item \textbf{Dữ liệu:} Sử dụng tập dữ liệu chuẩn LISA Traffic Sign Dataset.
        \item \textbf{Đánh giá phần cứng:} Các tổ hợp mô hình trên được đưa vào huấn luyện và chạy thử nghiệm thực tế trên nhiều kiến trúc phần cứng khác nhau bao gồm: CPU, GPU, TPU và Hệ thống nhúng (Embedded System).
        \item \textbf{Công cụ đo lường:} Đề xuất một phương pháp đánh giá dựa trên thống kê mô tả để đo lường hai yếu tố cốt lõi: thời gian huấn luyện/xử lý và độ chính xác nhận diện.
    \end{itemize}

    \subheading{Ưu điểm}
    \begin{itemize}
        \item Cung cấp một khung đánh giá có tính thực tiễn rất cao, đặc biệt hữu ích cho việc tối ưu hóa và đưa các mô hình AI xuống chạy trực tiếp trên các thiết bị biên (Edge/Embedded Devices) của xe tự hành.
        \item Phân tích chéo minh bạch giữa nhiều loại bộ vi xử lý, giúp dễ dàng xác định được "nút thắt cổ chai" (bottlenecks) về hiệu năng trong toàn bộ hệ thống.
        \item Sử dụng tập dữ liệu chuẩn hóa giúp các so sánh mang tính tin cậy cao.
    \end{itemize}

    \subheading{Nhược điểm}
    \begin{itemize}
        \item Các kiến trúc được chọn để đánh giá (MobileNet v1, ResNet v1, SSD) có phần hơi cũ so với các mô hình State-of-the-Art hiện nay (như các phiên bản YOLO mới hay DETR).
        \item Thử nghiệm chủ yếu tập trung trên tập dữ liệu LISA. Hiệu năng có thể thay đổi nếu áp dụng trên các tập dữ liệu có tính chất địa lý hoặc điều kiện môi trường nhiễu nặng hơn.
    \end{itemize}

    \subheading{Kết quả}
    \begin{itemize}
        \item Đánh giá thực nghiệm chứng minh rằng việc sử dụng TPU (Tensor Processing Unit) cho phép đạt tốc độ huấn luyện nhanh hơn tới 16,3 lần so với sử dụng GPU truyền thống.
        \item Ngoài lợi thế về tốc độ, TPU cũng mang lại kết quả tốt hơn về mặt độ chính xác (precision detection) cho một số tổ hợp mô hình cụ thể được thử nghiệm trong bài.
    \end{itemize}
\end{paperbox}

\vspace{0.5cm}

% --- BÀI SỐ 8 ---
\begin{paperbox}{8. A Comprehensive Survey and Analysis of Traffic Sign Recognition Systems With Hardware Implementation}
    \textbf{Nguồn:} \url{https://ieeexplore.ieee.org/document/10679144}\\

    \subheading{Vấn đề giải quyết}
    \begin{itemize}
        \item Sự phát triển của xe tự hành đòi hỏi các hệ thống Hỗ trợ Lái xe Tiên tiến (ADAS) ngày càng tinh vi, trong đó Nhận diện Biển báo Giao thông (Traffic Sign Recognition - TSR) đóng vai trò then chốt.
        \item Bài báo giải quyết nhu cầu cần có một cái nhìn tổng quan, hệ thống hóa lại toàn bộ các phương pháp cấu thành nên một hệ thống TSR (từ phát hiện đến phân loại). Đồng thời, giải quyết thách thức lớn trong thực tiễn: làm thế nào để triển khai các mô hình Deep Learning trong nhận diện biển báo giao thông lên các thiết bị phần cứng nhúng một cách hiệu quả để chạy được trong thời gian thực (real-time).
    \end{itemize}

    \subheading{Phương pháp}
    \begin{itemize}
        \item \textbf{Khảo sát toàn diện:} Phân tích và tổng hợp các kỹ thuật nhận diện đa dạng, từ các phương pháp truyền thống (dựa trên màu sắc, hình dạng) đến các thuật toán tiên tiến của Machine Learning và Deep Learning.
        \item \textbf{Thực nghiệm phần cứng:} Triển khai thử nghiệm trực tiếp các hệ thống TSR trên hai nền tảng phần cứng phổ biến là Raspberry Pi và Jetson Nano.
        \item \textbf{Đánh giá:} Đo lường và phân tích hiệu năng hoạt động của các thuật toán trên hai nền tảng này để tìm ra kiến trúc phần cứng phù hợp nhất.
    \end{itemize}

    \subheading{Ưu điểm}
    \begin{itemize}
        \item Cung cấp một bức tranh toàn cảnh rất cập nhật (năm 2024) về sự tiến hóa của hệ thống TSR, cực kỳ hữu ích cho việc xây dựng phần Tổng quan tài liệu (Literature Review) trong các dự án nghiên cứu khoa học.
        \item Không chỉ dừng lại ở lý thuyết thuật toán mà chú trọng mạnh mẽ vào tính ứng dụng thực tiễn thông qua việc đánh giá trực tiếp trên các thiết bị biên (edge devices).
        \item Đề xuất được chiến lược kết hợp đồng bộ (synergistic combination) giữa các phương pháp tốt nhất.
    \end{itemize}

    \subheading{Nhược điểm}
    \begin{itemize}
        \item Mang đặc thù của một bài báo khảo sát (Survey) kết hợp thực nghiệm, bài viết chủ yếu đánh giá lại các phương pháp đã có thay vì đề xuất một kiến trúc mạng Deep Learning hoàn toàn mới.
        \item Phạm vi thử nghiệm phần cứng chỉ giới hạn ở Raspberry Pi và Jetson Nano. Việc thiếu vắng các thử nghiệm trên các vi điều khiển hoặc phần cứng chuyên dụng thế hệ mới hơn có thể hạn chế một phần tính toàn diện trong bối cảnh công nghệ thay đổi nhanh chóng.
    \end{itemize}

    \subheading{Kết quả}
    \begin{itemize}
        \item Nghiên cứu cung cấp hệ thống dữ liệu chi tiết so sánh hiệu năng của các thuật toán nhận diện khi chạy trên Raspberry Pi và Jetson Nano, chỉ ra đâu là phần cứng tối ưu cho từng tình huống.
        \item Bài báo đưa ra kết luận và đề xuất xây dựng một phương pháp luận TSR thời gian thực mới bằng cách kết hợp hai phương pháp tốt nhất (về mặt phát hiện và phân loại). Sự kết hợp này hướng tới việc tối ưu hóa điểm cân bằng giữa độ chính xác (accuracy) và thời gian xử lý (processing time) khi triển khai lên nền tảng phần cứng đích.
    \end{itemize}
\end{paperbox}

\end{document}
